\runninghead{While Men Slept}
\properlane{word-in-the-heart}{While Men Slept: The Field Kept by the Word That Dwells in the Heart}

The enemy sows \latin{cum dormírent hómines}, ``while men were asleep,'' and
goes away. The parable does not say who the sleepers were or blame them, and
when the Lord explains the parable to his disciples he names the sower, the
field, the seed, the enemy, the harvest and the reapers, and passes over the
sleepers. The Fathers who comment on the words take them as a warning: the
enemy works while those who should keep the field are asleep, and St John
Chrysostom says that this kind of sleep can be avoided. In the Epistle Paul
asks the Colossians to let the word of Christ dwell in them abundantly, to let the peace of Christ
rejoice in their hearts, and to sing to God in their hearts. He says nothing
of sleep. Chrysostom, who warns of the sleepers, also makes the indwelling
word the reading of the Scriptures at home (\work{Hom.\ in Col.}\ 9), but he
does not set the two passages side by side. They are set together here
because both ask something of each hearer: the parable, as Chrysostom reads
it, asks him to stay awake, and Paul says what he is to keep in his heart.

\subsection{Who were asleep?}
\label{sec:heart-sleepers}

The words raise two questions the parable leaves open: whose sleep let the
enemy in, and was the sleep a fault? Sleep at night is no sin, and the
parable says the enemy came while ``men,'' not ``the servants,'' were asleep.

St Jerome gives the sleepers an office: \latin{Homines qui dormiunt,
magistros Ecclesiarum intellige}, ``by the men who sleep understand the
teachers of the Churches,'' and draws a rule for them: \latin{non dormiat, qui
Ecclesiae praepositus est, ne per illius negligentiam inimicus homo
superseminet zizania, hoc est, haereticorum dogmata}, ``let him who is set
over the Church not sleep, lest through his negligence the enemy sow cockle,
that is, the teachings of heretics'' (\work{In Matthaeum} II, on 13:36--43,
PL~26 col.~93).

St John Chrysostom also reads the words as a warning to those who keep the
field, and goes further: ``It is no small
danger, which He hereby suspends over our rulers, to whom especially is
entrusted the keeping of the field; and not the rulers only, but the subjects
too.'' He compares this sleep to the wayside, the rock and the thorns of the
sower: ``so here again sleep occasions our ruin; so that there is need of
continual watchfulness. Wherefore He also said, `He that endureth to the end,
the same shall be saved.''' Then he answers the obvious objection: ``And how
is it possible not to sleep? one may say. Indeed, as to natural sleep, it is
not possible; but as to that of our moral faculty, it is possible. Wherefore
Paul also said, `Watch ye, stand fast in the faith''' (\work{Hom.\ in Mt.}\
46.1).

Chrysostom applies the Gradual's psalm to the enemy the Lord names as the
sower of the cockle, the devil (Mt 13:39). At the verse before the Gradual's,
where the speakers will not trust in their own bow or sword, he tells his
hearer to say the same when he fights the devil: that he trusts in the mercy
of God and not in his own strength (\work{Expositio in Ps.}\ 43, §6). The
Gradual begins at the next verse, with thanks for a deliverance that did not
come by the speakers' bow: \latin{Liberásti nos, Dómine, ex affligéntibus
nos}, ``thou hast delivered us, O Lord, from them that afflict us.'' The watch
Chrysostom asks for at the parable is kept against the same enemy, so by his
own counsel the one who keeps it does not rely on his own strength.

The Collect asks God to keep his family \latin{contínua pietáte}, with
continual mercy, and to defend it always by his protection, because it leans
on the hope of heavenly grace alone. Since natural sleep, as Chrysostom says,
cannot be avoided, those who keep the field need a keeper themselves, and the
Collect asks God to be one.

The Lord's words that Chrysostom quotes promise salvation to the one who
endures to the end, and the Postcommunion asks for the effect of that
salvation, \latin{illíus salutáris \dots\ efféctum}, whose pledge has already
been received \latin{per hæc mystéria}.

St Rabanus Maurus applies the parable to each person. The good seed is sown
by the one \latin{qui in corpore suo semen bonae voluntatis ad bonorum
operum studet perducere fructum}, ``who strives in his own body to bring the
seed of good will to the fruit of good works.'' The sleepers warn us,
\latin{ut si semen bonum in agro nostro germinare velimus, inimici semina
caveamus. Ne dum somno inertiae torpemus, fructum laboris nostri perdamus}:
``if we want the good seed to sprout in our field, we must beware the enemy's
seed, lest while we grow torpid in the sleep of sloth we lose the fruit of
our labour'' (\work{In Matthaeum} IV, on 13:24--30, PL~107 col.~947).

St Thomas Aquinas notes how the enemy works: \latin{occultum se reddidit}, he
made himself hidden, and he sows \latin{in medio tritici, et fidelis populi},
in the midst of the wheat and of the faithful people (\work{Super Mt.}\
c.~13). He also reports Augustine's reading of \latin{cum crevisset herba},
``when the blade was sprung up,'' of the growth by which a man becomes
spiritual and able to discern (1~Cor 2:15).

Chrysostom extends the warning to every member of the Church, and Rabanus
puts the field in each person: it is ``our field,'' in which each sows the
seed of good will in his own body and must beware the enemy's seed. The Lord's
gloss makes the field the world and the good seed ``the children of the
kingdom'' (Mt 13:38), and Aquinas holds to it here. In the parable of the
sower, he says, the seed was the word of God sown in a man; in this parable
\latin{accipitur ipse homo, in quo seminatur}, the seed is the man himself in
whom it is sown, \latin{unde non est alia facienda expositio ab ea quam
Dominus fecit}, ``so no other explanation is to be given than the one the
Lord gave'' (\work{Super Mt.}\ c.~13). Since the seed is the man, the word
Paul asks to dwell in him is not the seed of this parable, and Chrysostom's
homily on Paul's verse gives it another use.

\subsection{The word that dwells abundantly}
\label{sec:heart-word}

Paul's verb for the word of Christ is \latin{hábitet}: let it dwell, take up
residence. St John Chrysostom stresses both words of the phrase. ``He
said not, `Let the word of Christ' be in you, simply, but what? `dwell in
you,' and `richly.''' He applies it at once to the households of his hearers:
``Hearken ye, as many as are worldly, and have the charge of wife and
children; how to you too he commits especially the reading of the Scriptures
and that not to be done lightly, nor in any sort of way, but with much
earnestness.'' He tells them to get books, at least the New Testament, and to
keep them as a chest of medicines, and he names the cause of the trouble:
``This is the cause of all evils, the not knowing the Scriptures. We go into
battle without arms, and how ought we to come off safe?'' (\work{Hom.\ in
Col.}\ 9).

Chrysostom does not connect this homily with the one on the sleepers, where
he asks for ``continual watchfulness.'' In both he describes an enemy at
work against people with no defence ready, and in this homily he names the
defence, the Scriptures read at home; so the sleep he warns of at the parable
has its remedy in the word he urges here.

St Thomas Aquinas makes Chrysostom's two points in the order of the verse. The
word of Christ is the origin of wisdom: \latin{ex verbo Christi hauriatis
eam}, draw wisdom from the word of Christ. Some do not have the word at all,
\latin{ideo nec sapientiam}, and so have no wisdom; therefore Paul says
\latin{Habitet}. Others are content with a little of it, \latin{sed Apostolus
vult quod habeamus multum}, ``but the Apostle wants us to have much,'' and
therefore he says \latin{abundánter}. The wisdom drawn from the word has
three uses, \latin{Triplex autem est usus hujus sapientiae, scilicet
instructionis, devotionis et directionis}: it instructs, it moves to
devotion, and it directs action (\work{Super Col.}\ c.~3 lect.~3).

Aquinas finds the third use of this wisdom, \latin{directionis in opere},
direction in work, in verse~17, ``all whatsoever you do in word or in work,
do all in the name of the Lord Jesus.'' The Secret asks God himself for that
direction: \latin{nutántia corda tu dírigas}, ``direct our wavering hearts,''
hearts that sway like a thing not firmly set.

\subsection{The peace that rules the heart: umpire and joy}
\label{sec:heart-peace}

At verse~15 the Missal sings \latin{pax Christi exsúltet in córdibus
vestris}, ``let the peace of Christ rejoice in your hearts.'' The Greek text
that St John Chrysostom expounds has a different verb, which means ``let it
act as umpire,'' and his text reads ``the peace of God.'' Chrysostom builds
his exposition on the umpire. ``If two thoughts are fighting together, set
not anger, set not spitefulness to hold the prize, but peace; for instance,
suppose one to have been insulted unjustly; of the insult are born two
thoughts, the one bidding him to revenge, the other to endure; and these
wrestle with one another: if the Peace of God stand forward as umpire, it
bestows the prize on that which bids endure, and puts the other to shame.
How? by persuading him that God is Peace, that He hath made peace with us.''
He sums it up: ``He hath represented an arena within, in the thoughts, and a
contest, and a wrestling, and an umpire'' (\work{Hom.\ in Col.}\ 8).
Theodoret, who reads the same Greek, has the peace in the heart preside over
the contest and award the prize, and work the concord dear to God
(\work{Interpretatio in Col.}, on 3:15). Theophylact follows Chrysostom's
picture: the peace of God stands as just arbiter of the inner contest and
gives the palm to the thought that bids long-suffering (\work{Expositio in
Col.}, on 3:15, PG~124 cols.~1261--1264).

St Thomas Aquinas expounds the Latin word the Mass sings. Peace springs from
charity, and \latin{Exultet, quia charitatis effectus est gaudium, quod
sequitur ex pace}: ``Let it rejoice, because the effect of charity is joy,
which follows from peace.'' The peace is Christ's, the peace he made between
God and man (\work{Super Col.}\ c.~3 lect.~3).

Both place in the heart a peace that God or Christ has already made with
man, and Chrysostom's umpire wins by persuading the man of it. The Introit
gives God's own words about his intention toward his people, \latin{Ego
cógito cogitatiónes pacis, et non afflictiónis}, ``I think thoughts of peace,
and not of affliction''; a man who believes them believes what Chrysostom's
umpire would persuade him of, that God is for peace with him.

\subsection{Singing in the heart}
\label{sec:heart-singing}

Verse~16 asks for ``psalms, hymns and spiritual canticles, singing in grace
in your hearts to God.'' St John Chrysostom: ``Not simply with the mouth, he
means, but with heedfulness. For this is to `sing to God,' but that to the
air, for the voice is scattered without result.'' He says that a man can sing
to God in the market-place with no one hearing him, as Moses ``cried in
thought'' and God answered, ``Why criest thou unto Me?'' (Ex 14:15): ``For it
is not forbidden one even when walking to pray in his heart, and to dwell
above'' (\work{Hom.\ in Col.}\ 9). Theodoret writes on the same verse that the
singing is in the heart, ``not with the mouth only'' (\work{Interpretatio in
Col.}, on 3:16).

St Thomas Aquinas reads the verse as the second use of this wisdom, devotion.
The canticles are spiritual \latin{quia quicquid nos facimus debemus referre
ad bona spiritualia, ad promissa aeterna, et ad reverentiam Dei}, ``because
whatever we do we must refer to spiritual goods, to the eternal promises, and
to the reverence of God''; and they are sung \latin{In cordibus, et non in
labiis tantum}, ``in the heart, and not on the lips only,'' for Isaias
complains of a people who honour God with their lips while their heart is far
from him (Is 29:13). Then he says what the Church's own singing is for:
\latin{Sunt autem cantica ecclesiae cordis principaliter, sed oris sunt ut
excitetur canticum cordis, et pro simplicibus et rudibus}. ``The songs of the
Church are principally of the heart; but they are of the mouth so that the
song of the heart may be roused, and for the sake of the simple and
unlettered'' (\work{Super Col.}\ c.~3 lect.~3). Aquinas says this of the
Church's songs in general, and it holds of the three psalms this Mass sings.

\subsection{The cry from the depths of the heart}
\label{sec:heart-depths}

The Alleluia and the Offertory sing \latin{De profúndis clamávi}. St Hilary
of Poitiers says that the prophet cries \latin{ex interno cordis sui secreto},
``from the inner secret of his heart,'' and tells the reader how to picture
him: \latin{existima illum non labiis, sed secreto, ut jussum est, cordis
orare}, ``think of him as praying not with his lips but in the secret of the
heart, as we are commanded'' (Mt 6:6) (\work{Tractatus in Ps.}\ 129.2, PL~9
cols.~719--720). The cry fits, he says, \latin{aut ignorationi, aut passioni, aut
orationi}: one who does not know, one who suffers, or one who prays.

St John Chrysostom says that \latin{de profundis} does not mean simply from
the mouth and tongue, since words can pour out while the mind wanders, but
from the deepest heart and the foundations of the mind. Prayer rooted so, like
a tree, stands against the devil's assaults; prayer from the lips alone
cannot rise to God (\work{Expositio in Ps.}\ 129, §1). St Robert
Bellarmine speaks of ``two depths, the depth of misery and the depth of his
heart,'' and of the second as ``a thorough consideration and deep reflection
on his own misery; for he that is not aware of, and that does not reflect on
the depth in which he lies, has no wish to rise out of it'' (\work{Explanatio
in Ps.}\ 129, vv.~1--2). St Augustine speaks of those who ``are very deep in
the deep, who do not even cry from the deep'' (\work{Enarr.\ in Ps.}\ 129.1).

Hilary, Chrysostom and Bellarmine speak of prayer at Psalm~129, not of the
Mass's singing, but all three put the cry in the heart, so the \latin{De
profúndis} this Mass sings three times is a singing ``in your hearts'' of the
kind Paul asks for. The man Bellarmine and Augustine describe, who does not
see the deep he lies in and so does not cry from it, fits Chrysostom's account
of the sleeper, asleep in his moral faculty.

The Communion promises an answer to prayer: \latin{Amen dico vobis, quidquid orántes
pétitis, crédite quia accipiétis, et fiet vobis}. Its first words open
verse~23 of Mark, where the one whose word casts a mountain into the sea is
the one who ``shall not stagger in his heart, but believe.'' St Bede reads
that mountain as the devil, named a mountain for the pride by which he lifts
himself against God, and says that it is taken up and cast into the sea at
the command of those strong in faith \latin{cum praedicantibus verbum
doctoribus sanctis immundus Spiritus ab eorum corde repellitur, qui ad vitam
sunt praeordinati}, ``when, as holy teachers preach the word, the unclean
spirit is driven out of the hearts of those ordained to life''; one who has
perfect faith, he adds, can move such mountains \latin{orando, vel etiam
jubendo}, by praying or even by commanding (\work{In Marci evangelium
expositio} III, on 11:23--25). Bede's devil is the parable's enemy, so his
reading has the sower of the cockle driven from the heart by the word preached
and by the faith that prays. The Communion makes its promise to that faith,
to whoever asks believing, and so to those who keep the field by prayer from
the heart, as Hilary and Chrysostom describe it.

\begin{fourSenses}
\item[Literal.] An enemy sows darnel among the wheat while men sleep, and it
appears only when the crop bears fruit; Paul tells the Colossians to let the
peace of Christ rejoice in their hearts, to let the word of Christ dwell in
them abundantly, and to sing to God in their hearts.
\item[Allegorical.] The Son of man sows the children of the kingdom in the
world and the devil sows his own among them (Mt 13:37--39, to which Aquinas
holds); the word of Christ is to dwell in each member of the Church richly
(Chrysostom, Aquinas); the devil is driven from the hearts of those ordained
to life when the word is preached (Bede); the peace in the heart is the peace
of God, ``that He hath made peace with us'' (Chrysostom), the peace Christ
made between God and man (Aquinas); and the Church's songs are sung by the
mouth to rouse the song of the heart (Aquinas).
\item[Moral.] Stay awake in the will, since natural sleep cannot be avoided
and moral sleep can (Chrysostom); guard the seed of good will against the
enemy's seed, lest the fruit of labour be lost in the torpor of sloth
(Rabanus); keep the Scriptures in the home (Chrysostom); let
the peace of God settle the inner contest between revenge and endurance
(Chrysostom); sing and pray with the heart and not the lips only (Chrysostom,
Aquinas, Hilary).
\item[Anagogical.] The watch is kept to the end, because ``He that endureth
to the end, the same shall be saved'' (Chrysostom on the sleepers, \work{Hom.\
in Mt.}\ 46.1); the spiritual canticles sung in the heart are referred
\latin{ad promissa aeterna}, to the eternal promises (Aquinas, \work{Super
Col.}\ c.~3 lect.~3); and the wheat, the children of the kingdom, is
gathered at the harvest, ``the wheat gather ye into my barn'' (Mt 13:30).
\end{fourSenses}
